\section{Benchmarking AMD MI300X for RLVR Training Workloads}

\smallskip\noindent\textit{Study Guide companion: Study Guide Section 8 (Research Proposals) and Section 12 (Hardware Deep Dive) cover the interview-ready version of this material.}
\smallskip

\begin{whythis}
Frontier labs currently run primarily on NVIDIA hardware. But RLVR training runs last \textit{weeks}, making every hardware decision a multi-billion-dollar call. AMD's MI300X offers 5.3~TB/s of HBM3 bandwidth (vs.\ H100's 3.35~TB/s) and 192~GB of on-chip memory (vs.\ 80~GB) --- advantages that map directly onto the memory-bandwidth-bound operations that dominate RLVR wall time. The engineering question is whether AMD's software ecosystem is mature enough for the hardware advantage to survive contact with real workloads. This chapter builds the evidence base. Simon Boehm did this for NVIDIA; nobody has done it rigorously for AMD RLVR workloads. That gap is wide open.
\end{whythis}

\subsection{Prerequisites and Hardware}

\textbf{Reading:} HIP Programming Guide (\texttt{rocm.docs.amd.com}); AMD CDNA3 Architecture Whitepaper (covers XCD topology, LDS, MFMA); \texttt{composable\_kernel} library documentation (AMD's equivalent of CUTLASS); \texttt{rocprof} and \texttt{omniperf} profiling guides.

\textbf{Hardware:} Rent MI300X on Crusoe or TensorWave at $\sim$\$1.70/hr. Budget $\sim$\$300/month for this chapter. All findings extend to the MI355X --- same CDNA architecture family, higher clocks.

\textbf{Setup:} ROCm 6.x stack, PyTorch with ROCm backend, \texttt{hipcc} compiler, \texttt{rocblas} and \texttt{hipblaslt} for BLAS baselines.

% -----------------------------------------------------------------------
\subsection{Part A: Conceptual Foundation}
% -----------------------------------------------------------------------

\begin{thinkfirst}{AMD vs.\ NVIDIA Architecture: What Changes About Your Mental Model?}
Before writing a single line of HIP, reason through the architectural differences. This shapes every tiling and scheduling decision.

\textbf{Execution width:} NVIDIA warps are 32 threads wide. AMD wavefronts are 64 threads wide. If you tile a GEMM with a 32-wide thread block strategy from a CUDA tutorial, you are leaving half the hardware idle on MI300X. How must your tile dimensions change? What happens to register pressure when you double the execution width?

\textbf{Local Data Share vs.\ shared memory:} MI300X LDS is 64~KB per compute unit; NVIDIA's shared memory is 48~KB (configurable up to 99~KB on H100 with \texttt{cudaFuncSetAttribute}). The MI300X has more on-chip scratch by default. For a double-buffered matmul tile that holds two BF16 sub-matrices of size $T \times T$: what is the largest $T$ you can hold in LDS before spilling? How does this compare to the H100 default?

\textbf{MFMA vs.\ Tensor Cores:} AMD MFMA (Matrix Fused Multiply-Add) instructions operate on 16$\times$16 or 32$\times$32 tiles. NVIDIA's \texttt{wmma}/MMA operates on 16$\times$16 tiles. Different instruction shapes mean different register-file layouts for accumulator tiles. Before writing Kernel~2, sketch out the data layout your wavefront needs in registers to feed a 32$\times$32 MFMA correctly.

\textbf{GCD topology:} MI300X has 8 Graphics Compute Dies (XCDs), each with its own LDS. Inter-XCD communication goes through the Infinity Fabric, not shared L1. A kernel that touches memory resident on a different XCD pays a higher penalty than expected. How would you detect this with \texttt{omniperf}? What workload patterns expose it?

\textbf{HBM3 bandwidth:} MI300X delivers 5.3~TB/s; H100 delivers 3.35~TB/s --- a 1.58$\times$ advantage. This advantage only materializes for memory-bandwidth-bound operations. Write down: which operations in an RLVR training loop do you expect to be memory-bandwidth-bound, and which compute-bound? Predict before you measure.
\end{thinkfirst}

\begin{thinkfirst}{Roofline Analysis for MI300X}
The roofline model places operations on a two-regime map: memory-bound (performance $\propto$ bandwidth) vs.\ compute-bound (performance $\propto$ FLOP/s peak). The crossover is the \textit{arithmetic intensity ridge point}.

\textbf{MI300X specs:} Peak BF16 compute $\approx 383$~TFLOP/s; peak memory bandwidth $= 5.3$~TB/s. Derive the ridge point:
$$I^* = \frac{383 \times 10^{12} \text{ FLOP/s}}{5.3 \times 10^{12} \text{ B/s}} \approx 72 \text{ FLOP/byte}$$

\textbf{Compare to H100:} Peak BF16 $\approx 989$~TFLOP/s (with sparsity; $\sim 495$ dense); bandwidth $= 3.35$~TB/s. H100 ridge point $\approx 148$~FLOP/byte (dense). The MI300X's ridge is lower --- it becomes compute-bound at \textit{lower} arithmetic intensity than H100. Counterintuitive result: for operations with intermediate arithmetic intensity ($72$--$148$~FLOP/byte), MI300X may deliver better absolute throughput even at lower peak FLOP/s.

\textbf{BF16 matmul arithmetic intensity:} An $N \times N$ GEMM ($C = AB$, $A,B,C$ in BF16) performs $2N^3$ FLOPs and reads/writes $3 \times 2N^2$ bytes. Arithmetic intensity $= \frac{2N^3}{6N^2} = N/3$. At what $N$ is MI300X compute-bound? ($N/3 > 72 \Rightarrow N > 216$.) For a 7B model, typical sequence batch matmuls have $N \gg 216$ --- firmly compute-bound during training.

\textbf{Autoregressive generation:} During decode, batch size is often 1--4 tokens. The weight matrix is $O(D^2)$ bytes; FLOPs are $O(D^2)$ per token but memory reads are also $O(D^2)$ --- arithmetic intensity $\approx 1$--$2$~FLOP/byte. Far below the ridge: \textit{pure memory-bandwidth bottleneck}. MI300X's 1.58$\times$ bandwidth advantage should give a near-linear generation speedup. Verify this in Part~B.

\textbf{Prediction log:} Before running a single benchmark, fill in this table for 5 RLVR operations: (operation, expected intensity, expected bound, expected AMD advantage). You will check this against measured results in Part~C.
\end{thinkfirst}

\begin{thinkfirst}{Why Generation Throughput Dominates RLVR Wall Time}
RLVR training iterates: generate rollouts with the current policy, score them with a verifier, compute policy gradient, update weights. Repeat for weeks.

\textbf{Amdahl's Law applied:} Suppose generation consumes 60\% of wall time and the training step (forward + backward + optimizer) consumes 40\%. If you optimize generation by 2$\times$:
$$\text{Speedup} = \frac{1}{0.40 + 0.60/2} = \frac{1}{0.70} \approx 1.43\times$$
If instead you optimize the matmul inside the training step by 2$\times$, and matmul is 50\% of that 40\% step (20\% of total wall time):
$$\text{Speedup} = \frac{1}{0.80 + 0.20/2} = \frac{1}{0.90} \approx 1.11\times$$

A 2$\times$ generation improvement beats a 2$\times$ matmul improvement by 29 percentage points of pipeline speedup. Over a 3-week run, the 1.43$\times$ speedup saves 8.5 days vs.\ 2.3 days.

\textbf{Why MI300X's bandwidth advantage is structural here:} Generation is memory-bandwidth-bound by the KV~cache and weight loading. No algorithm change fixes this --- the roofline ceiling is bandwidth. MI300X's 1.58$\times$ bandwidth advantage translates almost directly to 1.58$\times$ generation throughput. This is the hardware bet. Test whether the software stack delivers it.

\textbf{Dario's ``laminar flow'' principle:} Numerical stability during RL training matters. Unlike supervised learning, a gradient explosion mid-run wastes weeks of compute, not hours. Before optimizing for throughput, understand what precision regime keeps gradients stable. This is why Kernel~3 and Part~B Op~4 focus on FP8/BF16 mixed precision with gradient norm monitoring --- not just raw TFLOP/s.

\textbf{Write your answer before proceeding:} Which single operation should you optimize first on MI300X to maximize RLVR training throughput? Why? If your answer is not ``autoregressive generation,'' re-read this box.
\end{thinkfirst}

% -----------------------------------------------------------------------
\subsection{Part A: HIP Matmul Progression}
% -----------------------------------------------------------------------

\begin{exercise}{HIP Matmul: Naive to Near-rocBLAS in 4 Kernels (8--12 hours)}
This progression builds kernel credibility and MI300X proficiency. The pedagogy is in the profiling and diagnosis loop, not the kernel count. After each kernel: measure, profile with \texttt{omniperf}, form a hypothesis about the bottleneck, then write the next kernel to address exactly that bottleneck. Do not move on until you can articulate why each kernel is faster than the previous one.

\textbf{Kernel~0: rocBLAS Baseline.}
Call \texttt{rocblas\_sgemm} (or \texttt{hipblaslt} for BF16) for matrix sizes $N \in \{512, 1024, 2048, 4096, 8192\}$. Record TFLOP/s as your ceiling. Record bandwidth utilization. Compute achieved arithmetic intensity and verify it matches the roofline prediction from the Think First box. This number is your denominator for all subsequent kernels. If you cannot reach 80\% of theoretical peak with rocBLAS, something is wrong with your setup --- fix it before proceeding.

\textbf{Record:} TFLOP/s at each $N$, \% of BF16 theoretical peak (383 TFLOP/s), HBM bandwidth utilization.

\medskip
\textbf{Kernel~1: Naive HIP $\to$ LDS Tiling (Boehm stages 1--3 collapsed).}
Start with the simplest possible kernel: one thread computes one output element $C[i,j]$ by iterating over the $k$ dimension. Measure TFLOP/s --- expect $<5\%$ of rocBLAS. Then add memory coalescing (ensure threads in a wavefront access consecutive memory addresses) and LDS blocking in one progression. Tile size: start with $64 \times 64$; AMD-specific --- LDS is 64~KB, so a $64 \times 64$ BF16 tile is $64^2 \times 2 = 8$~KB; you can hold 8 such tiles. Double-buffer: load tile $k$ into LDS while computing tile $k-1$.

\textbf{AMD-specific pitfall:} Wavefronts are 64-wide. If your thread block is $32 \times 32 = 1024$ threads, you have 16 wavefronts. Ensure your LDS access pattern avoids bank conflicts for 64-wide accesses --- the bank structure is 32 banks, but the stride pattern differs from CUDA. Profile with \texttt{omniperf --dispatch-ids} to catch LDS bank conflicts directly.

\textbf{Record:} TFLOP/s, \% of rocBLAS, LDS bandwidth utilization, bank conflict rate.

\medskip
\textbf{Kernel~2: MFMA + Wavefront Tiling.}
This is the AMD-differentiating stage. Replace scalar FMA with MFMA instructions. Use \texttt{\_\_builtin\_amdgcn\_mfma\_f32\_32x32x8\_bf16} for 32$\times$32$\times$8 BF16 MFMA. Each MFMA instruction consumes registers from the full 64-wide wavefront --- you must lay out your register file so that each lane of the wavefront holds the correct $A$ and $B$ fragments. Tile dimensions: outer tile 128$\times$128 (computed by the thread block), inner tile 32$\times$32 (computed by one MFMA per wavefront). Load $A$ fragments (32$\times$8 BF16) and $B$ fragments (8$\times$32 BF16) from LDS into registers. Accumulate into 32$\times$32 FP32 accumulators.

\textbf{Profiling checkpoint:} Run \texttt{omniperf} and check VALU (vector ALU) utilization vs.\ MFMA utilization. If VALU is high and MFMA is low, your kernel is spending time on address computation, not matrix math. Reduce address arithmetic in the hot loop.

\textbf{Record:} TFLOP/s, \% of rocBLAS, VALU utilization, MFMA utilization, wavefront occupancy.

\medskip
\textbf{Kernel~3: FP8/BF16 Mixed Precision + Vectorized 128-bit Loads.}
RL training stability requires careful precision management. Use BF16 for weight storage and FP8 for activation matmul (where supported by ROCm), accumulating into FP32. This is the regime Dario's ``laminar flow'' principle applies to: monitor gradient norm variance across training steps. If norms spike, FP8 accumulation is the suspect --- fall back to BF16.

Replace scalar LDS loads with 128-bit vectorized loads (\texttt{float4} for 4$\times$FP32 or \texttt{short8} for 8$\times$BF16). This reduces the number of load instructions per LDS access and improves memory-level parallelism. AMD's memory subsystem prefers 128-bit aligned accesses --- verify alignment or profile for misalignment penalties.

\textbf{Target:} $>50\%$ of rocBLAS TFLOP/s by this kernel. If you fall short, profile with \texttt{omniperf} and identify the binding constraint before declaring done. Common culprits: register spilling (check SGPRs/VGPRs in \texttt{rocprof} output), LDS bank conflicts, insufficient wavefront occupancy.

\textbf{Record:} TFLOP/s, \% of rocBLAS, gradient norm variance (run 100 steps of a toy training loop with these matmuls), effective HBM bandwidth.

\medskip
\textbf{Diagnostic discipline:} After each kernel, write one paragraph: what was the bottleneck, how did you diagnose it, what did you change, what was the result. These paragraphs become the core of your blog post's technical narrative.

\textit{This establishes kernel credibility. The actual value for the portfolio is in Part~B.}
\end{exercise}

% -----------------------------------------------------------------------
\subsection{Part B: RLVR-Relevant Operations Benchmark}
% -----------------------------------------------------------------------

\begin{exercise}{RLVR Operations Benchmark on MI300X (15--20 hours)}
Part~A proved you can write HIP kernels. Part~B answers the question that matters: \textit{does MI300X's hardware advantage survive real RLVR workloads?} For each operation below: (1) state why it appears in RLVR training, (2) run the benchmark, (3) compare to H100 published numbers or your own H100 measurements if available.

\textbf{Benchmark harness:} Write a single Python script that runs all 7 operations, logs results to a JSON file, and prints a summary table. Use \texttt{torch.cuda.synchronize()} (which maps to \texttt{torch.hip.synchronize()} on ROCm) for accurate timing. Warm up 10 iterations; measure 100 iterations; report median and p95 latency.

\medskip
\textbf{Operation 1: Autoregressive Generation Throughput.} \textit{The RLVR bottleneck.}

RLVR spends 50--70\% of wall time generating rollouts. Each rollout is autoregressive: one token at a time, loading the full weight matrix and KV~cache at every step. This is memory-bandwidth-bound.

Load a 7B-parameter model in BF16 (\textasciitilde{}14~GB). Generate 512 tokens with batch sizes $\{1, 4, 8, 16, 32\}$. Measure tokens/second end-to-end (including KV~cache updates). MI300X's 5.3~TB/s should give $\approx 1.58\times$ vs.\ H100's 3.35~TB/s, since the bottleneck is weight loading bandwidth. Verify. If you observe less than $1.3\times$ speedup, profile KV~cache overhead separately --- the gap is likely in attention kernel efficiency or RCCL initialization overhead, not raw HBM bandwidth.

\textbf{Measure:} Tokens/second at each batch size. Roofline utilization (achieved HBM bandwidth / 5.3~TB/s). Ratio vs.\ H100 baseline.

\medskip
\textbf{Operation 2: Flash Attention Kernel.} \textit{The training compute kernel.}

Flash Attention fuses the attention computation to avoid materializing the full $N \times N$ attention matrix. On AMD, the reference implementation is in \texttt{composable\_kernel} (AMD's equivalent of CUTLASS). Alternatively, use \texttt{flash-attention} with the ROCm backend if available for your ROCm version.

Run for sequence lengths $\{512, 1024, 2048, 4096, 8192\}$, head dim 128, 32 heads, batch size 4. Compare to FlashAttention-2 on H100 (TFLOP/s numbers from Tri Dao's paper are public). Measure both forward and backward pass. If \texttt{composable\_kernel} FA underperforms by $>30\%$ vs.\ H100, identify whether the gap is in the MFMA utilization or the softmax numerics --- these have different fixes.

\textbf{Measure:} TFLOP/s (fwd), TFLOP/s (bwd), \% of theoretical MI300X peak, ratio vs.\ FlashAttention-2 H100.

\medskip
\textbf{Operation 3: All-Reduce (RCCL).} \textit{The distributed training bottleneck.}

Multi-GPU RLVR training synchronizes gradients via all-reduce after each backward pass. On NVIDIA, NCCL is battle-hardened. AMD uses RCCL (ROCm Collective Communication Library).

Using 4 MI300X GPUs (rent a 4-GPU node), run RCCL all-reduce at message sizes $\{1\text{MB}, 10\text{MB}, 100\text{MB}, 1\text{GB}, 10\text{GB}\}$. Plot bandwidth utilization (achieved GB/s / theoretical NVLink or Infinity Fabric bandwidth). Compare to published NCCL numbers at equivalent message sizes. Key question: does RCCL's bandwidth utilization fall off at small message sizes (the latency-dominated regime) --- and what is the crossover message size where it becomes bandwidth-efficient? For a 7B model, gradient tensors are $\sim$14~GB total --- is this in the efficient regime?

\textbf{Measure:} All-reduce bandwidth (GB/s) at each message size, \% of Infinity Fabric peak, ratio vs.\ NCCL H100.

\medskip
\textbf{Operation 4: Mixed-Precision Training Step.} \textit{Throughput + stability.}

A realistic training step: FP8 matmul for forward pass, BF16 for attention, FP32 accumulation for gradient updates. This matches the precision hierarchy used in production RLVR.

Construct a minimal training loop: 2-layer transformer, 7B-parameter size, batch size 4, sequence length 512. Run 100 forward-backward-optimizer steps. Measure: (a) TFLOP/s of the matmul-heavy forward pass, (b) gradient norm at each step --- flag if norm exceeds $10\times$ the initial norm (instability signal), (c) wall time per step.

This is numerically interesting because RL rewards are often sparse and bursty --- gradient norms spike around reward events. FP8 accumulation may be fine for supervised training but unstable for RL. Document what you find.

\textbf{Measure:} TFLOP/s, gradient norm variance, wall time/step, any instability events.

\medskip
\textbf{Operation 5: KV Cache Capacity.} \textit{The memory-capacity advantage.}

MI300X has 192~GB HBM3 vs.\ H100's 80~GB. For RLVR, long-context generation is increasingly important --- reasoning traces can be thousands of tokens. The question is at what context length MI300X's capacity advantage becomes the binding constraint.

For a 7B model in BF16: KV~cache per token per layer $= 2 \times 2 \times D_{\text{head}} \times H_{\text{heads}} = 2 \times 2 \times 128 \times 32 = 16,384$~bytes $= 16$~KB. For 32 layers: $512$~KB per token. At what context length does the KV~cache alone fill 80~GB? (Answer: $\approx 160K$ tokens.) At what length does it fill 192~GB? ($\approx 375K$ tokens.) Verify these numbers experimentally by running generation until OOM. Repeat for a 70B model.

\textbf{Measure:} Max context length before OOM for 7B and 70B models. Tokens/second as context length grows (bandwidth demand grows with KV~cache size). The point at which throughput degrades due to KV~cache bandwidth pressure.

\medskip
\textbf{Operation 6: MoE Memory Budget.} \textit{Dario's MoE bet.}

Frontier labs increasingly use MoE architectures. An 8-expert MoE (2 active per token) has $8\times$ the total parameter count but only $2\times$ the active FLOPs. The memory footprint for all experts must fit on-device.

Model: 8 experts, each a 7B MLP block (2-layer FFN, hidden dim 28,672). Expert weights: $8 \times 7B \times 2~\text{bytes} = 112$~GB for BF16. On H100 (80~GB): does not fit. On MI300X (192~GB): fits with 80~GB remaining for KV~cache, activations, and optimizer state. At what context length does the remaining 80~GB for KV~cache run out? (Use Op~5 calculation for 7B equivalent.) At what total model scale does MI300X's 192~GB advantage become the reason you pick AMD over NVIDIA?

\textbf{Measure:} Expert weight memory footprint. Remaining capacity for KV~cache. Max context length achievable with full MoE model loaded. Tokens/second for MoE decode vs.\ dense decode (expert routing overhead).

\medskip
\textbf{Operation 7: Decode-Path Attention (KV Cache Memory Movement).} \textit{Not matmul.}

Autoregressive decode is not a matmul problem --- it is a KV~cache memory movement problem. During decode step $t$, the model reads the full KV~cache (all $t$ previous tokens), computes attention scores, and writes one new KV pair. This is an irregular memory access pattern that MFMA instructions do not directly accelerate.

Use \texttt{composable\_kernel} decode attention kernels (or implement via \texttt{xformers} with ROCm backend). Profile with \texttt{omniperf}: what fraction of the decode step is (a) KV~cache reads, (b) attention score computation, (c) KV write. At sequence length 4096, what is the ratio of memory bandwidth consumed to FLOPs performed? Does this match the $\approx 1$~FLOP/byte prediction from the Think First roofline box?

\textbf{Measure:} KV~cache read bandwidth (GB/s), FLOPs/s for attention scoring, HBM bandwidth utilization as \% of 5.3~TB/s, ratio vs.\ H100 decode attention benchmark.
\end{exercise}

% -----------------------------------------------------------------------
\subsection{Part C: Intelligence-per-Watt Comparison}
% -----------------------------------------------------------------------

\begin{thinkfirst}{Predict Before You Measure: Intelligence-per-Watt}
Before running Part~C benchmarks, write down your predictions. Return to this box and grade yourself after.

MI300X TDP: 750~W. H100 SXM TDP: 700~W. The MI300X consumes 7\% more power. For AMD to win on intelligence-per-Watt, it must deliver $>7\%$ more useful throughput on a given operation.

\textbf{Predict for each Part~B operation:} Will MI300X deliver better FLOP/s-per-Watt (or tokens/s-per-Watt) than H100? Reason from first principles using the roofline analysis and architecture differences you studied.

\begin{enumerate}
  \item Autoregressive generation: AMD wins/loses? Why? (Hint: bandwidth is proportional to power consumption approximately linearly for HBM.)
  \item Flash Attention forward: AMD wins/loses? What determines the efficiency ratio?
  \item All-reduce (RCCL vs.\ NCCL): AMD wins/loses? What is the primary source of software overhead?
  \item Mixed-precision training step: AMD wins/loses? Does FP8 throughput on MFMA vs.\ Tensor Cores change the picture?
  \item KV~cache capacity (192~GB vs.\ 80~GB): This is not a throughput question --- it is a capability question. At what context length does MI300X unlock workloads that are simply impossible on H100?
\end{enumerate}

\textbf{Write your predictions now.} Do not read Part~C exercise results first.
\end{thinkfirst}

\begin{exercise}{Intelligence-per-Watt Analysis (2--3 hours)}
Using the benchmark results from Parts~A and~B, build the decision-useful analysis that justifies or refutes AMD investment.

\textbf{Step 1: Compute useful FLOP/s per Watt.} For each Part~B operation, record: (achieved throughput on MI300X) / 750~W and (H100 benchmark throughput) / 700~W. Use tokens/s for generation; TFLOP/s for matmul and attention; GB/s for all-reduce. Build a table with columns: Operation, MI300X metric, H100 metric, MI300X/W, H100/W, Ratio.

\textbf{Step 2: Gap analysis.} For operations where MI300X loses on per-Watt efficiency, categorize the gap source:
\begin{itemize}
  \item \textit{Software immaturity} --- kernel not yet optimized (e.g., RCCL, FA not tuned): gap closable with engineering effort.
  \item \textit{Compiler gap} --- \texttt{hipcc} generating worse code than \texttt{nvcc} for a class of patterns: harder to close, may require upstream ROCm contributions.
  \item \textit{Architectural mismatch} --- the operation genuinely suits NVIDIA's architecture better (e.g., Tensor Core scheduling, warp-level primitives): structural disadvantage, cannot be engineered away cheaply.
\end{itemize}
For each gap, estimate engineering effort to close: $<1$~person-week, 1--4~person-weeks, $>1$~person-month, or ``architectural, cannot close.''

\textbf{Step 3: Write the memo.} One page, structured as follows:
\begin{enumerate}
  \item \textbf{Recommendation} (one sentence): Should Anthropic invest in MI300X for RLVR training, and under what conditions?
  \item \textbf{Hardware case} (2--3 bullets): The specific operations where AMD wins on intelligence-per-Watt and why the win is structural.
  \item \textbf{Risks} (2--3 bullets): The specific gaps, their sources, and engineering cost to close.
  \item \textbf{Decision trigger} (1 paragraph): What would need to be true for the answer to change? E.g., ``If RCCL reaches NCCL parity at the 10~GB message size, the all-reduce gap closes and the hardware case strengthens by X\%.'' This is the section a senior RL engineer can act on.
\end{enumerate}

\textbf{Verification:} Your analysis should identify at least one operation where AMD wins on intelligence-per-Watt AND at least one where it loses. If your analysis shows AMD wins everywhere, you have not stress-tested the software gaps. If it shows AMD loses everywhere, you have not accounted for the bandwidth and memory-capacity advantages.
\end{exercise}

% -----------------------------------------------------------------------
\subsection{Portfolio Project}
% -----------------------------------------------------------------------

\begin{portfolio}{Evaluating AMD MI300X for Large-Scale RLVR Training}
\textbf{Deliverable:} Public blog post + open-source repository. Estimated total time: 30--40 hours including Parts A--C.

\textbf{Repository structure:}
\begin{itemize}
  \item \texttt{kernels/} --- HIP kernels from Part~A, each in a separate file with a build script.
  \item \texttt{benchmarks/} --- Python benchmark harness for all 7 Part~B operations, runnable on any MI300X with ROCm 6.x.
  \item \texttt{results/} --- Raw JSON output from your benchmark runs, plus a Jupyter notebook that generates all plots.
  \item \texttt{analysis/} --- The intelligence-per-Watt memo and gap analysis from Part~C.
  \item \texttt{README.md} --- Reproduction instructions: exact ROCm version, model checkpoint, cloud provider, instance type, and total cost to reproduce.
\end{itemize}

\textbf{Blog post structure:} Part~A kernels are the pedagogical hook that draws readers who want to learn HIP. Parts~B and~C are the decision-useful substance that RL infrastructure teams can act on. The narrative thread: ``I set out to understand whether AMD's hardware advantage survives real RLVR workloads. Here is what I found, operation by operation.''

\textbf{The monkey:} Does MI300X's 1.58$\times$ bandwidth advantage actually translate to faster RLVR training end-to-end, or does software ecosystem immaturity (RCCL, composable\_kernel, compiler maturity) erase the hardware advantage? The benchmark harness answers this empirically. The gap analysis answers what would need to change for the answer to flip. That combination --- empirical data plus actionable gap analysis --- is what makes the project decision-useful rather than merely educational.

\textbf{Verification criteria:}
\begin{enumerate}
  \item The analysis identifies $\geq 1$ operation where AMD wins on intelligence-per-Watt.
  \item The analysis identifies $\geq 1$ operation where AMD loses, with a categorized gap source.
  \item The gap analysis includes a concrete engineering effort estimate ($<$1 week / 1--4 weeks / $>$1 month / architectural) for each gap.
  \item The benchmark harness is reproducible by a third party with access to a single MI300X instance at the stated ROCm version.
  \item Total cloud cost to reproduce is documented and is under \$200.
\end{enumerate}

\textbf{Why this matters:} Simon Boehm's CUDA matmul post demonstrated kernel authorship on NVIDIA. This project does the same for AMD, but at a higher level of strategic relevance: it produces evidence that informs a hardware purchasing decision for any team evaluating AMD for RL workloads. A senior RL engineer reading your blog can directly use the gap analysis and cost estimates. That is the difference between a tutorial and a contribution.
\end{portfolio}

\newpage
